% Accepted B.87: separate Part B input, reviewed after checkpoint 15.
\statement{B.87}{Proved}{The actual source is exactly the fixed subgroup on the both-family filled surface, discretely and in full completion}
\label{stmt:B87}
Let
\[
 Y_b=\{(v,w)\in(\mathbb C^*\setminus\{1\})^2:
           v\ne w,\ vw\notin\mu _5\},\qquad
 L_b=Y_b\cap\{v+w=2\},
\]
with the actual base point $P=(399/200,1/200)$. The rational involution
\[
 F(v,w)=\left(\frac{v(w-1)}{1-v},
                    \frac{w(v-1)}{1-w}\right)
\]
preserves $Y_b$ and fixes $P$ and the whole source $L_b$ pointwise.
Put $\Pi_b=\pi _1(Y_b,P)$, $\sigma=F_*$, and
$H_b=\operatorname{im}(\pi _1(L_b,P)\to\Pi_b)$. Then
\begin{equation}\label{eq:B87-fixed}
 C_{\Pi_b}(\sigma)=H_b\cong F_{11},\qquad
 C_{\widehat\Pi_b}(\widehat\sigma)=\overline{H_b}
                                       \cong\widehat F_{11}.
\end{equation}
Here the bar is the actual closed image in the full completion. Its full
induced topology is part of the proof. The original source $L$ before the
constant and moving fillings has the same image in $\Pi_b$ and its full
completion. If $\tau$ denotes the specified external involution in
$\Pi_b\rtimes_\sigma C_2$, the equivalent centralizer statements are
\[
 C_{\Pi_b\rtimes C_2}(\tau)=H_b\times\langle\tau\rangle,
 \qquad
 C_{\widehat\Pi_b\rtimes C_2}(\tau)
                  =\overline{H_b}\times\langle\tau\rangle.
\]
This identifies the fixed subgroup after both families have been filled.
It does not identify the fixed subgroup before the symmetric deletions
are filled, or assert fixed-element lifting through their intersection
kernel.

\begin{proof}
\emph{B.87.1. The actual product bundle and fixed-point cover.}
Use the product coordinate $q=vw$, distinct from the tangential coordinate
of B.85. Its base is $B=\mathbb C^*\setminus\mu _5$, and its fiber is
\[
 X_q=\mathbb P^1\setminus
              \{0,\infty,1,q,+\sqrt q,-\sqrt q\}.
\]
The six points are distinct because $q\ne0,1$. On a sufficiently small
simply connected base neighborhood they can be labeled continuously;
isotopy extension of this finite moving set supplies actual local
trivializations. Thus $Y_b\to B$ is a six-punctured-sphere bundle.
The base and fiber are aspherical, with groups $\mathcal B=F_6$ and
$N=F_5$. Its homotopy exact sequence, using $\pi _2(B)=0$, gives
\begin{equation}\label{eq:B87-extension}
 1\longrightarrow N\longrightarrow\Pi_b
       \xrightarrow{p_*}\mathcal B\longrightarrow1.
\end{equation}
It splits as a group extension by choosing lifts of the six free base
generators. This noncanonical group splitting does not assert a global
geometric section of the bundle.

On the $q$-fiber the given involution is
\[
 F_q(v)=\frac{v-q}{v-1},\qquad
 \begin{pmatrix}1&-q\\1&-1\end{pmatrix}^{\!2}=(1-q)I.
\]
It exchanges the pairs $\{0,q\}$, $\{1,\infty\}$, and
$\{\sqrt q,-\sqrt q\}$. Its two fixed points are
$v=1\pm\sqrt{1-q}$, exactly the equation $v+w=2$ on the fiber.
Consequently its complete fixed locus in $Y_b$ is $L_b$.

The actual source projection is the unramified double cover
\[
 L_b\longrightarrow B,\qquad w(2-w)=q,
                    \qquad (w-1)^2=1-q.
\]
Both possible branch values $1,\infty$ are removed from $B$.
At $q_*=399/40000$ the chosen source point has
$v=399/200=1+\sqrt{1-q_*}$. Let
\[
 \epsilon:\mathcal B\longrightarrow C_2
\]
record the sign of $\sqrt{1-q}$. It sends the base meridian at $q=1$
to the nonidentity element and the meridians at $0$ and the four other
fifth roots to the identity. Its kernel $\mathcal B_0$ has index two
and free rank $11$. The source group maps isomorphically to
$\mathcal B_0$, hence injects into $\Pi_b$. Its actual image supplies
\begin{equation}\label{eq:B87-source-split}
 \Pi_0:=p_*^{-1}(\mathcal B_0)=N\rtimes H_b,
               \qquad H_b\xrightarrow{\sim}\mathcal B_0.
\end{equation}
Because $F$ fixes every source point, $\sigma$ fixes every element of
$H_b$ literally.

\emph{B.87.2. The fiber orbifold and the finite cone-class detector.}
Put $c=1-q$, $u=v-1$. The fiber involution becomes $u\mapsto c/u$,
with quotient coordinate $a=u+c/u$. Its three deleted quotient values
are $\infty,-2,q-2$, and its two order-two cone points are
$\pm2\sqrt{1-q}$. The deleted values are pairwise distinct, the cone
points are distinct, and none of the cone points is deleted: the possible
coincidences would force $q=0$ or $q=1$.

Write $\Gamma=\Pi_b\rtimes\langle\tau\rangle$, where conjugation
by $\tau$ is $\sigma$. Since $\sigma$ acts trivially on the base,
the kernel of $\Gamma\to\mathcal B$ is the actual fiber orbifold group
\[
 E=N\rtimes\langle\tau\rangle\cong C_2*C_2*F_2.
\]
Indeed a sphere with three punctures and two order-two cone points has
presentation
\[
 \langle a_0,a_1,d_0,d_1,d_\infty:
 a_0^2=a_1^2=1,\ a_0a_1d_0d_1d_\infty=1\rangle;
\]
eliminating $d_\infty$ gives the displayed free product. Its actual
index-two manifold cover is $X_q$. The specified based involution
$\tau$ represents the cone at the chosen fixed point; it may be taken
as $a_0$ in this presentation.

The two cone classes are transported exactly as the two geometric fixed
points. To see the marking explicitly, on a small base neighborhood
choose the two roots of $1-q$, trivialize the moving quotient punctures
and cone points, and lift the isotopy through the double orbifold cover.
These are local equivariant trivializations of the actual fiber, so the
permutation of cone classes along a base loop is precisely $\epsilon$.
No unmarked identification is substituted for this fixed-point cover.

In particular this permutation is detected in the characteristic finite
quotient
\begin{equation}\label{eq:B87-cone-vectors}
 A=E/([E,E]E^2)\cong\mathbb F_2^4.
\end{equation}
The two cone elements have distinct independent images $\overline a_0$
and $\overline a_1$. Conjugation by $\Pi_b$ permutes these two vectors
by $\epsilon\circ p_*$. Conjugation by $N$ is inner in $E$ and acts
trivially on $A$. This finite detector will also exclude additional
outside-cover fixed elements in the full completion.

\emph{B.87.3. The discrete fixed subgroup.}
The reduced-word normal form for $C_2*C_2*F_2$ gives
$C_E(\tau)=\langle\tau\rangle$ and hence $C_N(\sigma)=1$.
If $g\in\Pi_b$ is fixed by $\sigma$, its conjugation fixes
$\overline a_0$ in \eqref{eq:B87-cone-vectors}. The cone permutation
forces $p_*(g)\in\mathcal B_0$. Write $g=nh$ using
\eqref{eq:B87-source-split}, with $n\in N$ and $h\in H_b$.
Since $h$ is fixed pointwise, $\sigma(g)=g$ implies $\sigma(n)=n$,
so $n=1$. The converse is the literal pointwise source invariance.
This proves the first equality of \eqref{eq:B87-fixed} without assuming
that every fixed group element already has a pointwise fixed loop
representative.

\emph{B.87.4. Actual fiber topology and exact completion.}
We verify the completion step directly. Finite-index characteristic
subgroups are cofinal in the finitely generated free group $N$: intersect
all subgroups of bounded index. For one such subgroup $K$, the action of
the chosen splitting $\mathcal B\to\Pi_b$ on $N/K$ has finite image
$A_K\leq\operatorname{Aut}(N/K)$. Therefore
\[
 \Pi_b\longrightarrow (N/K)\rtimes A_K
\]
is a finite quotient whose kernel intersects $N$ in exactly $K$.
Thus $N$ carries its full profinite topology from $\Pi_b$.
Elements outside $N$ are separated by finite quotients of $\mathcal B$;
in particular the extension is residually finite.

The closure of $N$ is consequently $\widehat N$. The finite quotients
of $\Pi_b$ killing $N$ are exactly the finite quotients of $\mathcal B$,
so
\[
 1\longrightarrow\widehat N\longrightarrow\widehat\Pi_b
       \longrightarrow\widehat{\mathcal B}\longrightarrow1
\]
is exact. The chosen splitting extends continuously. Since
$\mathcal B_0$ has finite index, its completion is its actual closure
in $\widehat{\mathcal B}$. The source isomorphism
$H_b\simeq\mathcal B_0$ therefore gives an injective map on full
completions, with image $\overline{H_b}$, and
\begin{equation}\label{eq:B87-completed-source-split}
 \widehat\Pi_0=\widehat N\rtimes\overline{H_b}
              =p_*^{-1}(\widehat{\mathcal B_0}).
\end{equation}
This is the actual open subgroup of $\widehat\Pi_b$, and the completed
$\sigma$ fixes $\overline{H_b}$ pointwise.

Completion of the finite extension $E=N\rtimes C_2$ is
$\widehat N\rtimes C_2$. The universal property of finite quotients
of a free product also gives
\[
 \widehat E=C_2\amalg C_2\amalg\widehat F_2.
\]
The full-profinite centralizer theorem of
\citet[Theorem B$'$]{A_HR1985} gives
$C_{\widehat E}(\tau)=\langle\tau\rangle$. Hence
$C_{\widehat N}(\widehat\sigma)=1$. This uses the full-category
finite-factor centralizer theorem, not a pro-$p$ recognition theorem or
a Kurosh assertion for arbitrary closed subgroups.

The finite quotient $A$ in \eqref{eq:B87-cone-vectors} remains the
mod-two abelianization of $\widehat E$. Conjugation by
$\widehat\Pi_b$ acts continuously on it and extends the actual finite
action of $\Pi_b$. Its permutation on the two cone vectors is therefore
the completed character $\widehat\epsilon\circ p_*$.
A completed element centralizing $\tau$ fixes $\overline a_0$ and must
belong to the open subgroup \eqref{eq:B87-completed-source-split}.
Writing it as $nh$ there, its fixedness and the fiber centralizer equality
force $n=1$. This proves the second equality of
\eqref{eq:B87-fixed}. The stated centralizers in the split extensions
follow immediately.

\emph{B.87.5. Original source and the two remaining symmetric-deletion gates.}
The original source $L$ is obtained from $L_b$ by deleting the four
$w=u$ and four $w=2-u$ points for $u\in\mu _5\setminus\{1\}$.
Filling them is surjective on fundamental groups and full completions.
Thus its actual images in the filled surface are precisely $H_b$ and
$\overline{H_b}$.

For the symmetric open subset $Y\cap F^{-1}(Y)\subset Y_b$, the
identification of its own fixed subgroup with its source image is a
separate question. Also separate is whether a fixed element in the
fiber product of the two one-sided filling quotients admits a fixed lift
through the intersection of their two actual kernels. Neither assertion
is inferred from \eqref{eq:B87-fixed}; in particular no whole marginal
or joint-source kernel equality follows from this statement alone.

\emph{Verification record.}
The actual rational involution, all three puncture exchanges, the quotient
coordinate, the literal base point, and the source-cover discriminant
were checked symbolically in \path{work_A/code/}, file
\path{check_filled_surface_involution.py}, with genuine saved output.
The cited full-profinite centralizer theorem and its proof were reopened
in the primary author-associated text, printed pages159--160. The
completion and finite cone-permutation arguments are proved above rather
than supplied by a generic exactness assertion.
\end{proof}
